% COLLECTION_METADATA: {"schema_version": 1, "kind": "reused_writeup", "independently_reviewed": false, "primary_source": "attacks/erdos/gpt_pro_5.4/256.tex", "source_paths": ["attacks/erdos/gpt_pro_5.2/256.tex", "attacks/erdos/gpt_pro_5.4/256.tex"], "source_model": "gpt_pro_5.4", "source_completion": 100.0, "source_claim": "solved"}
\section*{Erd\H{o}s Problem 256 --- collection record}

\subsection*{PROVENANCE AND STATUS}
Official problem: https://www.erdosproblems.com/256

Saved collaborative-database status: open.
Snapshot checked: 2026-09-23T17:22:24Z.
Latest status and formalization links: https://teorth.github.io/erdosproblems/

This is an attributed collection of existing writeups. The model-folder names below identify their original attribution; placement in GPT\_6\_Astra\_Ultra does not attribute their mathematical work to that model. All locally available source versions selected for this problem are reproduced below, without editing their text.

Proof claims, computations, literature statements, and completion estimates in the reused text have not been independently verified in this collection pass. A source claiming a solution does not change the saved upstream status.

Primary source (latest numbered version in the preferred model folder): \texttt{attacks/erdos/gpt\_pro\_5.4/256.tex}.

\subsection*{REUSED SOURCE 1}
Original attribution: \texttt{gpt\_pro\_5.2}.
Original file: \texttt{attacks/erdos/gpt\_pro\_5.2/256.tex}.

% BEGIN REUSED SOURCE: attacks/erdos/gpt_pro_5.2/256.tex
\section*{Erd\H{o}s Problem \#256}

\subsection*{1) FORMAL RESTATEMENT}
For each integer $n\ge 1$ and each nondecreasing $n$-tuple $(a_1,\dots,a_n)\in\mathbb{N}^n$ (i.e. $a_1\le\cdots\le a_n$), set
\[
P_{a_1,\dots,a_n}(z) \;:=\; \prod_{i=1}^n (1-z^{a_i}),\qquad z\in\mathbb{C}.
\]
For such a tuple define
\[
M(a_1,\dots,a_n)\;:=\;\max_{\left\lvert z\right\rvert=1}\,\left\lvert P_{a_1,\dots,a_n}(z)\right\rvert\in[0,\infty).
\]
(The maximum exists because $z\mapsto\left\lvert P_{a_1,\dots,a_n}(z)\right\rvert$ is continuous on the compact set $\{z\in\mathbb{C}:\left\lvert z\right\rvert=1\}$.)
Define
\[
f(n)\;:=\;\inf_{a_1\le\cdots\le a_n\in\mathbb{N}} M(a_1,\dots,a_n).
\]
Equivalently, $f(n)$ is the largest real number such that for \emph{every} choice of $a_1\le\cdots\le a_n\in\mathbb{N}$ one has
\[
\max_{\left\lvert z\right\rvert=1}\left\lvert \prod_{i=1}^n (1-z^{a_i})\right\rvert\;\ge\; f(n).
\]

The specific sub-question is the truth of the asymptotic lower bound
\[
(\star)\qquad \exists\,c>0\ \exists\,K>0\ \exists\,n_0\ \forall\,n\ge n_0:\ \log f(n)\ge K n^c.
\]
Here $\log$ denotes the natural logarithm; changing the base rescales $K$ by a constant factor.

\subsection*{2) QUICK LITERATURE/CONTEXT CHECK}
(From the standard summary on this problem and the cited papers.)
\begin{itemize}[leftmargin=2.2em]
\item Erd\H{o}s--Szekeres (1959) showed $f(n)>\sqrt{2n}$ and $\lim_{n\to\infty} f(n)^{1/n}=1$ \cite{ErSz59}.
\item A long sequence of upper bounds culminates in the result of Belov--Konyagin (1996):
\[\log f(n)\ll (\log n)^4\quad(n\to\infty),\]
which in particular gives a negative answer to $(\star)$ \cite{BeKo96}.
\item Very recently, Tang (2025) improved the classical general lower bound to
\[ f(n)\ge 2\sqrt{n}\quad\text{for all }n\in\mathbb{N},\]
the first improvement of the constant in the general lower bound since 1959 \cite{Tang25}.
\end{itemize}

\subsection*{3) ATTACK PLAN}
We target $(\star)$.  The plan is:
\begin{enumerate}[leftmargin=2.2em]
\item Use the known upper bound $\log f(n)\ll(\log n)^4$ (Belov--Konyagin).
\item Prove the elementary comparison $(\log n)^4 = o(n^c)$ for every fixed $c>0$.
\item Conclude that no fixed $c>0$ can satisfy $(\star)$.
\end{enumerate}

\subsection*{4) WORK}
\paragraph{Tiny cases / sanity checks.}
For $n=1$ and any $a_1\in\mathbb{N}$,
\[
M(a_1)=\max_{\left\lvert z\right\rvert=1}\left\lvert 1-z^{a_1}\right\rvert=2,
\]
because $\left\lvert 1-w\right\rvert\le 2$ for $\left\lvert w\right\rvert=1$ and equality holds at $w=-1$ (so take $z=e^{i\pi/a_1}$). Hence
\[
f(1)=2.
\]
More generally $f(n)\le 2^n$ by taking $a_1=\cdots=a_n=1$, which gives $(1-z)^n$ and $\max_{\left\lvert z\right\rvert=1}\left\lvert 1-z\right\rvert=2$.

As a numerical experiment (not a proof): restricting to $a_i\le 30$ and sampling $\theta$ on a fine grid for $z=e^{i\theta}$ suggests candidates giving
\[
M(4,8)\approx 3.0792\quad\text{(so }f(2)\le 3.0792\text{ in this restricted search)},
\]
\[
M(2,4,6)\approx 4.3899\quad\text{(so }f(3)\le 4.3899\text{ in this restricted search)}.
\]
These are only heuristic upper bounds for the true $f(n)$.

\paragraph{Known upper bound (input).}
We use the following result.
\begin{theorem}[Belov--Konyagin, consequence for $f(n)$ \cite{BeKo96}]
There exists an absolute constant $C>0$ such that for all $n\ge 2$,
\[
\log f(n)\le C\,(\log n)^4.
\]
Equivalently, $f(n)\le \exp\bigl(C(\log n)^4\bigr)$ for all $n\ge 2$.
\end{theorem}

\paragraph{Elementary comparison lemma.}
\begin{lemma}\label{lem:log4_vs_power}
For every $c>0$, one has $(\log n)^4 = o(n^c)$ as $n\to\infty$. Equivalently,
\[\lim_{n\to\infty}\frac{(\log n)^4}{n^c}=0.\]
\end{lemma}
\begin{proof}
Fix $c>0$ and set $\varepsilon:=c/8$. Consider the function $g(x)=x^{\varepsilon}/\log x$ for $x>e$.
Then
\[
\log g(x)=\varepsilon\log x-\log\log x\xrightarrow[x\to\infty]{}+\infty,
\]
so $g(x)\to\infty$. In particular, there exists $x_0$ such that for all $x\ge x_0$,
\[\log x \le x^{\varepsilon}.
\]
For $n\ge x_0$ this gives $(\log n)^4 \le n^{4\varepsilon}=n^{c/2}$, hence
\[
0\le \frac{(\log n)^4}{n^c}\le \frac{n^{c/2}}{n^c}=n^{-c/2}\xrightarrow[n\to\infty]{}0.
\]
\end{proof}

\paragraph{Disproof of $(\star)$.}
\begin{proposition}\label{prop:no_power_lower}
Statement $(\star)$ is false; i.e. there is \emph{no} constant $c>0$ such that $\log f(n)\gg n^c$.
\end{proposition}
\begin{proof}
Assume for contradiction that $(\star)$ holds. Then there exist $c>0$, $K>0$ and $n_0$ such that
$\log f(n)\ge K n^c$ for all $n\ge n_0$.

On the other hand, by the Belov--Konyagin bound there is $C>0$ such that for all $n\ge 2$,
$\log f(n)\le C(\log n)^4$.
By Lemma~\ref{lem:log4_vs_power}, $(\log n)^4=o(n^c)$, so we may choose $n_1$ such that for all $n\ge n_1$,
\[C(\log n)^4 < K n^c.
\]
Taking $n\ge \max\{n_0,n_1,2\}$ gives
\[\log f(n)\le C(\log n)^4 < K n^c \le \log f(n),
\]
a contradiction.
\end{proof}

\subsection*{5) VERIFICATION}
\begin{itemize}[leftmargin=2.2em]
\item The logical negation of $(\star)$ is: for every $c>0$ and every $K>0$ there exist arbitrarily large $n$ with $\log f(n)<K n^c$. Proposition~\ref{prop:no_power_lower} proves this, in fact giving $\log f(n)=O((\log n)^4)$ for all $n$.
\item Lemma~\ref{lem:log4_vs_power} is a standard calculus comparison and is proved explicitly.
\item The argument uses only: (i) the stated Belov--Konyagin upper bound, and (ii) the elementary fact that polylogarithms are $o(n^c)$.
\end{itemize}

\subsection*{6) FINAL}
\noindent LABEL: \textbf{FULL SOLUTION}\\
\noindent SUBLABEL: \textbf{COUNTEREXAMPLE/DISPROOF}\\
The claim ``$\exists c>0$ such that $\log f(n)\gg n^c$'' is false.  Indeed Belov--Konyagin prove
$\log f(n)\ll(\log n)^4$ \cite{BeKo96}, and since $(\log n)^4=o(n^c)$ for every $c>0$ (Lemma~\ref{lem:log4_vs_power}), no power lower bound $\log f(n)\gg n^c$ can hold.

\subsection*{7) COMPLETION ESTIMATE (MANDATORY)}
\noindent COMPLETION: 100\%.

%%%%%%%%%%%%%%%%%%%%%%%%%%%%%%%%%%%%%%%%%%%%%%%%%%%%%%%%%%%%%%%%%%%%%%%%%%%%%%%

% END REUSED SOURCE: attacks/erdos/gpt_pro_5.2/256.tex

\subsection*{REUSED SOURCE 2}
Original attribution: \texttt{gpt\_pro\_5.4}.
Original file: \texttt{attacks/erdos/gpt\_pro\_5.4/256.tex}.

% BEGIN REUSED SOURCE: attacks/erdos/gpt_pro_5.4/256.tex
\section*{Erd\H{o}s Problems \#256--\#258: Round 2 continuation}

\section*{Erd\H{o}s Problem \#256}

\subsection*{1) ROUND-2 OBJECTIVE}
Path \textbf{(C)}: obstruction/correction. Round 1 already gave a complete disproof of the specific subquestion
\[
\exists c>0\text{ such that }\log f(n)\gg n^c.
\]
So the only legitimate round-2 task here is to adversarially check that disproof, lock in the quantifiers, and record the minimal corrected conclusion.

\subsection*{2) Round-1 FOUNDATION USED}
I use the following Round-1 facts without re-proving them.
\begin{enumerate}
\item Belov--Konyagin's upper bound: there exists an absolute constant $C>0$ such that
\[
\log f(n)\le C(\log n)^4
\]
for all sufficiently large $n$.
\item The elementary comparison
\[
(\log n)^4=o(n^c)\qquad (n\to\infty)
\]
for every fixed $c>0$.
\item The Round-1 contradiction argument combining (1) and (2).
\end{enumerate}

\subsection*{3) NEW INSIGHT / TOOL (ROUND-2)}
There is no new mathematical ingredient needed. The only new point is an adversarial consistency check against the current Erd\H{o}s Problems summary, which now explicitly records that the specific power-lower-bound question has a negative answer.

\subsection*{4) ATTACK PLAN (ROUND-2)}
The only possible gap after Round 1 would have been a quantifier slip: one must rule out \emph{every} fixed exponent $c>0$, not just a particular one. The plan is therefore to re-run the contradiction with the quantifiers written explicitly.

\subsection*{5) WORK (ROUND-2)}
Fix any $c>0$. Suppose, for contradiction, that there exist $K>0$ and $n_0$ such that
\[
\log f(n)\ge K n^c\qquad (n\ge n_0).
\]
Round 1 already established that there is an absolute constant $C>0$ such that
\[
\log f(n)\le C(\log n)^4
\]
for all sufficiently large $n$. Since $(\log n)^4=o(n^c)$, there exists $n_1$ such that
\[
C(\log n)^4<K n^c\qquad (n\ge n_1).
\]
Hence, for all sufficiently large $n$,
\[
\log f(n)\le C(\log n)^4<K n^c\le \log f(n),
\]
which is impossible. Therefore no fixed $c>0$ can satisfy $\log f(n)\gg n^c$.

Thus the specific subquestion in Problem \#256 is fully settled in the negative.

\subsection*{6) ADVERSARIAL VERIFICATION}
\begin{enumerate}
\item \textbf{Quantifiers.} The argument starts with an \emph{arbitrary} fixed $c>0$ and shows that this $c$ cannot work. Hence it rules out all $c>0$.
\item \textbf{Dependence of constants.} The Belov--Konyagin constant $C$ is absolute; the contradiction only needs that $C$ is independent of $n$.
\item \textbf{Logical target.} The original question only asked whether a power-type lower bound $\log f(n)\gg n^c$ could hold. The argument disproves exactly that claim, not the entire surrounding estimation problem for $f(n)$.
\end{enumerate}
No gap remains.

\subsection*{7) FINAL}
\textbf{FULL SOLUTION --- COUNTEREXAMPLE / DISPROOF}

The specific claim
\[
\exists c>0\text{ such that }\log f(n)\gg n^c
\]
is false. Belov--Konyagin give the upper bound $\log f(n)\ll (\log n)^4$, and $(\log n)^4=o(n^c)$ for every fixed $c>0$. Hence no power lower bound of the form $\log f(n)\gg n^c$ can hold.

\subsection*{8) COMPLETION ESTIMATE (MANDATORY)}
COMPLETION: 100\%

\subsection*{9) REFERENCES}
\begin{enumerate}
\item A. S. Belov and S. V. Konyagin, \emph{An estimate of the free term of a non-negative trigonometric polynomial with integer coefficients}, Izv. Math. \textbf{60} (1996), 1123--1182.
\item T. F. Bloom (ed.), \emph{Erd\H{o}s Problem \#256}, \texttt{erdosproblems.com/256}, accessed 2026-03-07.
\end{enumerate}


% END REUSED SOURCE: attacks/erdos/gpt_pro_5.4/256.tex

\subsection*{COMPLETION ESTIMATE}
Inherited primary-source estimate: $100\%$. This is the original source's unverified estimate, not a new assessment or confirmation that the problem is solved.
